\section{Conductor module}
\label{sec:conductor}
% Coil-agnostic: given (B_peak, I_op, tau_dump, SC option) returns the
% cable. Same function is called by TF, CS, PF, mirror and (future)
% stellarator solvers. Code: cicc.m, heat_balance_cicc_ode.m, Ic_*.m,
% Jc_REBCO.m, size_cicc_cable.m.

\subsection{Superconductor selection and critical current}
Three options are available: pure LTS, pure HTS (REBCO) and hybrid. In the
LTS option NbTi is used below a threshold field and Nb$_3$Sn above; in the
hybrid option REBCO is used above a second threshold and LTS below
\todo{thresholds differ today: TF 6/15~T, CS--PF 6.2/14.5~T; make them one
input of the shared module}, with a fallback to REBCO when the LTS cable would
exceed the maximum number of strands. Critical-current parametrisations:
\begin{itemize}
\item Nb$_3$Sn: \todo{scaling law and strand data sets (DTT TF, ITER, WST)
-- cite source of each fit};
\item NbTi: \todo{scaling law};
\item REBCO: $I_c(B,T)$ fit at $T=20$~K with field angle dependence
\todo{cite fit origin}.
\end{itemize}
The number of superconducting strands/tapes is $N_\mathrm{sc}=\lceil
I_\mathrm{op}/I_c(B_\mathrm{peak})\rceil$ \todo{discuss operating margin:
temperature margin / I/Ic fraction -- currently implicit}.

\subsection{Quench protection: copper from hot-spot analysis}
The segregated copper is sized by integrating the adiabatic heat balance of
the cable during an exponential current dump,
\begin{equation}
  \sum_k A_k\,\rho_k c_{p,k}(T)\,\frac{\mathrm{d}T}{\mathrm{d}t} =
  \eta_\mathrm{Cu}(T,B)\,\frac{I(t)^2}{A_\mathrm{Cu}},\qquad
  I(t)=I_\mathrm{op}\,\mathrm{e}^{-(t-\tau_\mathrm{del})/\tau_\mathrm{dump}},
\end{equation}
\todo{check exact form against heat\_balance\_cicc\_ode.m: materials
included, helium, RRR, magnetoresistance.}
The number of copper strands is found by bisection until the hot-spot
temperature equals $T_\mathrm{lim}$ (250~K for LTS, 150~K for HTS).
The dump time constant is
$\tau_\mathrm{dump}=\max(\tau_\mathrm{VV},\,L I_\mathrm{op}/V_\mathrm{max},\,4~\mathrm{s})$.
The first term is the shortest discharge compatible with the eddy-current
loads on the vacuum vessel \cite{Giannini2023MADE,Bachmann_aspect},
\begin{equation}
\tau_\mathrm{VV}=\frac{B_0\,NI\,N_\mathrm{TF}\,(R_0/A)^2}{R_\mathrm{VV}\,\sigma_\mathrm{a,VV}},
\end{equation}
with $\sigma_\mathrm{a,VV}=120$~MPa; the second is set by the maximum
terminal voltage. \todo{justify the 4~s floor (hard-coded) or make it an
input.}

\subsection{Cable and jacket cross-section}
The cable area includes strands (void fraction, twist factor), the central
cooling channel and a 0.4~mm wrap. The jacket is sized for a given
conductor width $w_c$ (imposed by the winding-pack layout): for a
rectangular conductor the jacket thickness is increased until the radial
membrane stress in the innermost jacket layer,
\begin{equation}
  \sigma_\mathrm{rm} = p_\mathrm{m}\,\frac{w_c-2t_\mathrm{ins}}{2t_J}\,
  \frac{K_\mathrm{jacket}}{K_\mathrm{e,rad}} + \sigma_{z},\qquad
  p_\mathrm{m}=\frac{B_\mathrm{peak}^2}{2\mu_0},
\end{equation}
falls below $\sigma_\mathrm{allow}/s_m$. The equivalent radial and
toroidal stiffnesses $K_\mathrm{e,rad}$, $K_\mathrm{e,tor}$ of the
jacket--insulation--cable composite are obtained from series/parallel
spring models (Appendix~\ref{app:stiffness}). These smeared stiffnesses
are the same quantities that feed the beam sections of the global model
(section~\ref{sec:stellarator}).
